\documentclass[12pt,a4paper]{article}
\usepackage[utf8]{inputenc}
\usepackage{amsmath}
\usepackage{geometry}
\usepackage{hyperref}
\usepackage{listings}
\usepackage{xcolor}
\usepackage{tcolorbox}

\geometry{margin=2.5cm}

% Define colors for code highlighting
\definecolor{codebackground}{rgb}{0.95,0.95,0.95}
\definecolor{codekeyword}{rgb}{0.0,0.0,0.5}
\definecolor{codestring}{rgb}{0.6,0.1,0.1}
\definecolor{codecomment}{rgb}{0.25,0.5,0.35}

% Configure listings for shell commands
\lstdefinestyle{shellstyle}{
    backgroundcolor=\color{codebackground},
    basicstyle=\ttfamily\small,
    breaklines=true,
    captionpos=b,
    frame=single,
    keepspaces=true,
    showspaces=false,
    showstringspaces=false,
    showtabs=false,
    tabsize=2
}

\lstset{style=shellstyle}

\title{Terminal Commands Reference\\Speech Enhancement Project}
\author{}
\date{}

\begin{document}

\maketitle
\tableofcontents
\newpage

\section{Environment Setup}

\subsection{Create Virtual Environment}

\begin{lstlisting}[language=bash]
# Create virtual environment
python -m venv .venv

# Activate (Windows PowerShell)
.\.venv\Scripts\Activate.ps1

# Activate (Windows CMD)
.venv\Scripts\activate.bat

# Activate (Linux/Mac)
source .venv/bin/activate
\end{lstlisting}

\subsection{Install Dependencies}

\begin{lstlisting}[language=bash]
# Install all required packages
pip install -r requirements.txt

# Install in development mode (editable)
pip install -e .

# Verify installation
pip list
python -c "import torch; print(torch.__version__)"
\end{lstlisting}

\subsection{Check Hardware Configuration}

\begin{lstlisting}[language=bash]
# Check CUDA availability
python -c "import torch; print('CUDA:', torch.cuda.is_available())"

# Check CUDA version
python -c "import torch; print('CUDA Version:', torch.version.cuda)"

# Check GPU device
python -c "import torch; print('GPU:', torch.cuda.get_device_name(0))"
\end{lstlisting}

\section{Data Preparation}

\subsection{Generate Noisy Mixtures}

Mix clean LibriSpeech speech with DEMAND noise at multiple SNR levels:

\begin{lstlisting}[language=bash]
# Generate train set
python -m src.data_prep.mix_librispeech_with_demand --split train

# Generate validation set
python -m src.data_prep.mix_librispeech_with_demand --split val

# Generate test set
python -m src.data_prep.mix_librispeech_with_demand --split test

# Generate all splits sequentially
python -m src.data_prep.mix_librispeech_with_demand --split train
python -m src.data_prep.mix_librispeech_with_demand --split val
python -m src.data_prep.mix_librispeech_with_demand --split test
\end{lstlisting}

\textbf{Default SNR levels:} 0 dB, 5 dB, 10 dB

\subsection{Generate VAD Labels}

Generate ground truth Voice Activity Detection labels from clean speech:

\begin{lstlisting}[language=bash]
# Generate VAD labels for train set
python -m src.data_prep.generate_vad_labels --split train

# Generate VAD labels for validation set
python -m src.data_prep.generate_vad_labels --split val

# Generate VAD labels for test set
python -m src.data_prep.generate_vad_labels --split test

# Generate for all splits with custom parameters
python -m src.data_prep.generate_vad_labels \
    --split train \
    --energy-threshold 0.1 \
    --zcr-threshold 0.2
\end{lstlisting}

\textbf{Output:} Binary .npy files in \texttt{data/librispeech\_demand/vad\_labels/}

\subsection{Generate Dataset Metadata}

Extract comprehensive metadata about audio files:

\begin{lstlisting}[language=bash]
# Generate metadata for LibriSpeech-DEMAND
python -m src.data_prep.generate_metadata \
    --dataset librispeech \
    --output results/metadata_librispeech.json

# Generate metadata for VoiceBank-DEMAND
python -m src.data_prep.generate_metadata \
    --dataset voicebank \
    --output results/metadata_voicebank.json

# Generate for both datasets
python -m src.data_prep.generate_metadata --dataset librispeech
python -m src.data_prep.generate_metadata --dataset voicebank
\end{lstlisting}

\subsection{Validate Dataset}

Check dataset integrity and statistics:

\begin{lstlisting}[language=bash]
# Validate LibriSpeech-DEMAND dataset
python -m src.data_prep.validate_dataset \
    --dataset librispeech \
    --check-audio \
    --check-labels

# Validate VoiceBank-DEMAND dataset
python -m src.data_prep.validate_dataset \
    --dataset voicebank \
    --check-audio

# Quick validation (no audio loading)
python -m src.data_prep.validate_dataset --dataset librispeech

# Detailed validation with audio checks
python -m src.data_prep.validate_dataset \
    --dataset librispeech \
    --check-audio \
    --check-labels \
    --verbose
\end{lstlisting}

\section{Model Training}

\subsection{Train MaskNet Model}

Train the deep learning enhancement model:

\begin{lstlisting}[language=bash]
# Basic training with default parameters
python -m src.train.train_mask_model

# Training with custom parameters
python -m src.train.train_mask_model \
    --epochs 50 \
    --batch-size 8 \
    --lr 0.001 \
    --base-channels 32 \
    --use-lstm \
    --use-attention

# Training with reduced model (faster, less memory)
python -m src.train.train_mask_model \
    --base-channels 16 \
    --no-lstm \
    --no-attention \
    --batch-size 16

# Short test run (30 batches)
python -m src.train.train_mask_model \
    --max-batches 30 \
    --base-channels 16

# Resume training from checkpoint
python -m src.train.train_mask_model \
    --resume checkpoints/train_balanced_20260202_202648/masknet_best.pth

# Training with custom dataset split
python -m src.train.train_mask_model \
    --train-split train \
    --val-split val \
    --epochs 30

# Training with specific device
python -m src.train.train_mask_model \
    --device cuda \
    --batch-size 16

# Training without CUDA (CPU only)
python -m src.train.train_mask_model \
    --device cpu \
    --batch-size 4 \
    --num-workers 0
\end{lstlisting}

\subsection{Training Parameters}

\textbf{Architecture:}
\begin{itemize}
    \item \texttt{--base-channels}: Number of base channels (16, 32, 64)
    \item \texttt{--use-lstm}: Enable LSTM layers in bottleneck
    \item \texttt{--use-attention}: Enable channel attention
    \item \texttt{--no-lstm}: Disable LSTM (faster)
    \item \texttt{--no-attention}: Disable attention (faster)
\end{itemize}

\textbf{Training:}
\begin{itemize}
    \item \texttt{--epochs}: Number of training epochs (default: 50)
    \item \texttt{--batch-size}: Batch size (default: 8)
    \item \texttt{--lr}: Learning rate (default: 0.001)
    \item \texttt{--max-batches}: Stop after N batches (for testing)
    \item \texttt{--resume}: Resume from checkpoint path
\end{itemize}

\textbf{Hardware:}
\begin{itemize}
    \item \texttt{--device}: cuda or cpu
    \item \texttt{--num-workers}: DataLoader workers (0 for Windows)
    \item \texttt{--no-amp}: Disable mixed precision (FP16)
\end{itemize}

\subsection{Monitor Training}

\begin{lstlisting}[language=bash]
# View training logs in real-time
Get-Content checkpoints/train_*/training.log -Wait -Tail 50

# Check tensorboard logs
tensorboard --logdir checkpoints/train_*

# Check checkpoint info
python -c "import torch; ckpt = torch.load(
    'checkpoints/masknet_best.pth', 
    map_location='cpu', weights_only=False
); print('Epoch:', ckpt['epoch']); 
   print('Train Loss:', ckpt['metrics']['train_loss']); 
   print('Val Loss:', ckpt['metrics']['val_loss'])"

# List all checkpoints
Get-ChildItem checkpoints/train_*/*.pth | 
    Select-Object Name, Length, LastWriteTime
\end{lstlisting}

\section{Model Evaluation}

\subsection{Evaluate All Methods}

Compare all enhancement methods (MaskNet, Spectral Subtraction, Wiener Filter):

\begin{lstlisting}[language=bash]
# Evaluate with best checkpoint
python -m src.eval.evaluate_all_methods \
    --checkpoint checkpoints/masknet_best.pth \
    --output results/all_methods_results.json \
    --plot

# Evaluate on specific test set
python -m src.eval.evaluate_all_methods \
    --checkpoint checkpoints/masknet_best.pth \
    --test-dir data/librispeech_demand/test \
    --output results/evaluation.json

# Evaluate with custom parameters
python -m src.eval.evaluate_all_methods \
    --checkpoint checkpoints/train_*/checkpoint_epoch10_final.pth \
    --output results/epoch10_results.json \
    --max-files 100 \
    --plot

# Fast evaluation (10 files only)
python -m src.eval.evaluate_all_methods \
    --checkpoint checkpoints/masknet_best.pth \
    --max-files 10 \
    --no-plot
\end{lstlisting}

\textbf{Output files:}
\begin{itemize}
    \item JSON results: \texttt{results/all\_methods\_results.json}
    \item Comparison plots: \texttt{results/plots/method\_comparison.png}
    \item SNR plots: \texttt{results/plots/snr\_comparison.png}
\end{itemize}

\subsection{Evaluate VAD Baseline}

Evaluate Voice Activity Detection (Energy+ZCR method):

\begin{lstlisting}[language=bash]
# Evaluate on test set (all SNR levels)
python -m src.eval.evaluate_vad_baseline

# Evaluate specific SNR level
python -m src.eval.evaluate_vad_baseline --snr 5

# Evaluate with custom parameters
python -m src.eval.evaluate_vad_baseline \
    --snr 0 \
    --energy-threshold 0.1 \
    --zcr-threshold 0.2 \
    --output results/vad_snr0_custom.json

# Evaluate all SNR levels separately
python -m src.eval.evaluate_vad_baseline --snr 0
python -m src.eval.evaluate_vad_baseline --snr 5
python -m src.eval.evaluate_vad_baseline --snr 10

# Generate plots
python -m src.eval.evaluate_vad_baseline \
    --plot \
    --output results/vad_results.json
\end{lstlisting}

\textbf{Metrics computed:}
\begin{itemize}
    \item Accuracy, Precision, Recall, F1-Score
    \item Miss Rate, False Alarm Rate
\end{itemize}

\section{Visualization and Analysis}

\subsection{Plot Enhancement Metrics}

Generate comparison plots from evaluation results:

\begin{lstlisting}[language=bash]
# Plot results from JSON file
python -m src.plots.plot_enhancement_metrics \
    --input results/all_methods_results.json \
    --output results/plots/enhancement_comparison.png

# Plot with custom style
python -m src.plots.plot_enhancement_metrics \
    --input results/evaluation.json \
    --output results/plots/metrics.png \
    --dpi 300 \
    --figsize 12 8
\end{lstlisting}

\subsection{Plot Example Signals}

Visualize waveforms and spectrograms:

\begin{lstlisting}[language=bash]
# Plot single example
python -m src.plots.plot_example_signals \
    --clean data/librispeech_demand/test/clean/file.flac \
    --noisy data/librispeech_demand/test/noisy/file.wav \
    --enhanced demo_results/enhanced.wav \
    --output results/plots/signal_example.png

# Plot with spectrograms
python -m src.plots.plot_example_signals \
    --clean data/.../clean.flac \
    --noisy data/.../noisy.wav \
    --enhanced results/.../enhanced.wav \
    --show-spectrogram \
    --output results/plots/spectrogram_comparison.png
\end{lstlisting}

\subsection{Plot VAD Metrics}

Visualize VAD performance across SNR levels:

\begin{lstlisting}[language=bash]
# Plot VAD results
python -m src.plots.plot_vad_metrics \
    --input results/vad_baseline_metrics.json \
    --output results/plots/vad_performance.png

# Plot with custom configuration
python -m src.plots.plot_vad_metrics \
    --input results/vad_results.json \
    --output results/plots/vad_snr_analysis.png \
    --dpi 300
\end{lstlisting}

\section{Demo and Presentation}

\subsection{Run Complete Demo}

Interactive demonstration of all functionalities:

\begin{lstlisting}[language=bash]
# Run demo with best checkpoint
python demo.py --checkpoint checkpoints/masknet_best.pth

# Run demo with custom checkpoint
python demo.py \
    --checkpoint checkpoints/train_*/checkpoint_epoch20.pth

# Run demo on specific test file
python demo.py \
    --checkpoint checkpoints/masknet_best.pth \
    --test-dir data/voicebank_demand/test

# Run demo with custom output directory
python demo.py \
    --checkpoint checkpoints/masknet_best.pth \
    --output-dir demo_results_custom
\end{lstlisting}

\textbf{Demo includes:}
\begin{itemize}
    \item Voice Activity Detection visualization
    \item Spectral Subtraction enhancement
    \item Wiener Filter enhancement
    \item MaskNet enhancement
    \item Metric computation (PESQ, STOI, SNR)
    \item Waveform and spectrogram plots
    \item Audio file output for listening
\end{itemize}

\section{Utility Commands}

\subsection{Check Configuration}

\begin{lstlisting}[language=bash]
# View current configuration
python -c "from src import config; 
print('Sample Rate:', config.SAMPLE_RATE); 
print('FFT Size:', config.N_FFT); 
print('Hop Length:', config.HOP_LEN)"

# Check dataset paths
python -c "from src import config; 
print('LibriSpeech:', config.LIBRISPEECH_ROOT); 
print('VoiceBank:', config.VOICEBANK_ROOT)"

# Print all config values
python -m src.config
\end{lstlisting}

\subsection{Dataset Statistics}

\begin{lstlisting}[language=bash]
# Count files in dataset
python -c "from pathlib import Path; 
from src import config; 
train_files = list((config.LIBRISPEECH_ROOT / 
    'noisy' / 'train').glob('*.wav')); 
print(f'Training files: {len(train_files)}')"

# Check dataset size
Get-ChildItem data/librispeech_demand/noisy -Recurse | 
    Measure-Object -Property Length -Sum | 
    Select-Object Count, @{Name='SizeGB';
        Expression={[math]::Round($_.Sum / 1GB, 2)}}

# List checkpoint files
Get-ChildItem checkpoints -Recurse -Filter *.pth | 
    Format-Table Name, 
        @{Name='SizeMB';Expression={[math]::Round($_.Length/1MB,2)}},
        LastWriteTime
\end{lstlisting}

\subsection{Clean Up}

\begin{lstlisting}[language=bash]
# Remove Python cache
Get-ChildItem -Path . -Include __pycache__ -Recurse -Force | 
    Remove-Item -Force -Recurse

# Remove old checkpoints (keep best only)
Get-ChildItem checkpoints/train_*/checkpoint_epoch*.pth | 
    Where-Object {$_.Name -notmatch 'best'} | 
    Remove-Item

# Remove temporary files
Remove-Item -Recurse -Force demo_results/
Remove-Item -Recurse -Force logs/temp/

# Clean results directory
Remove-Item results/*.json
Remove-Item results/plots/*.png
\end{lstlisting}

\subsection{Export Results}

\begin{lstlisting}[language=bash]
# Create results archive
Compress-Archive -Path results/* -DestinationPath results.zip

# Create data archive (noisy only)
Compress-Archive -Path data/librispeech_demand/noisy/* 
    -DestinationPath data.zip

# Export checkpoint with metadata
python -c "import torch; import json; 
ckpt = torch.load('checkpoints/masknet_best.pth', 
    map_location='cpu', weights_only=False); 
metadata = {'epoch': ckpt['epoch'], 
    'metrics': ckpt['metrics'], 
    'config': ckpt['model_config']}; 
json.dump(metadata, open('checkpoint_info.json', 'w'), indent=2)"
\end{lstlisting}

\section{Troubleshooting}

\subsection{CUDA Out of Memory}

\begin{lstlisting}[language=bash]
# Reduce batch size
python -m src.train.train_mask_model --batch-size 4

# Reduce model size
python -m src.train.train_mask_model \
    --base-channels 16 \
    --no-lstm \
    --no-attention

# Disable mixed precision
python -m src.train.train_mask_model --no-amp

# Use CPU
python -m src.train.train_mask_model --device cpu
\end{lstlisting}

\subsection{DataLoader Errors (Windows)}

\begin{lstlisting}[language=bash]
# Set num_workers to 0
python -m src.train.train_mask_model --num-workers 0

# Verify dataset integrity first
python -m src.data_prep.validate_dataset --dataset librispeech
\end{lstlisting}

\subsection{Import Errors}

\begin{lstlisting}[language=bash]
# Reinstall dependencies
pip install --upgrade -r requirements.txt

# Install in editable mode
pip install -e .

# Check Python path
python -c "import sys; print('\n'.join(sys.path))"

# Verify imports
python -c "from src import config; from src.models import MaskNet"
\end{lstlisting}

\subsection{Audio File Errors}

\begin{lstlisting}[language=bash]
# Check audio file integrity
python -c "import soundfile as sf; 
audio, sr = sf.read('path/to/file.wav'); 
print(f'Shape: {audio.shape}, SR: {sr}')"

# Convert sample rate if needed
python -c "import soundfile as sf; 
import librosa; 
audio, _ = librosa.load('input.wav', sr=16000); 
sf.write('output.wav', audio, 16000)"

# Validate all audio files
python -m src.data_prep.validate_dataset \
    --dataset librispeech \
    --check-audio \
    --verbose
\end{lstlisting}

\section{Quick Reference}

\subsection{Complete Pipeline}

Run the entire pipeline from scratch:

\begin{lstlisting}[language=bash]
# 1. Setup environment
python -m venv .venv
.\.venv\Scripts\Activate.ps1
pip install -r requirements.txt

# 2. Prepare data (if not done)
python -m src.data_prep.mix_librispeech_with_demand --split train
python -m src.data_prep.mix_librispeech_with_demand --split val
python -m src.data_prep.mix_librispeech_with_demand --split test

# 3. Generate VAD labels
python -m src.data_prep.generate_vad_labels --split train
python -m src.data_prep.generate_vad_labels --split val
python -m src.data_prep.generate_vad_labels --split test

# 4. Validate dataset
python -m src.data_prep.validate_dataset --dataset librispeech

# 5. Train model
python -m src.train.train_mask_model \
    --epochs 50 \
    --batch-size 8 \
    --base-channels 32 \
    --use-lstm \
    --use-attention

# 6. Evaluate all methods
python -m src.eval.evaluate_all_methods \
    --checkpoint checkpoints/masknet_best.pth \
    --output results/all_methods_results.json \
    --plot

# 7. Evaluate VAD
python -m src.eval.evaluate_vad_baseline

# 8. Run demo
python demo.py --checkpoint checkpoints/masknet_best.pth
\end{lstlisting}

\subsection{Fast Testing Pipeline}

Quick test run for development:

\begin{lstlisting}[language=bash]
# 1. Short training (30 batches)
python -m src.train.train_mask_model \
    --max-batches 30 \
    --base-channels 16

# 2. Quick evaluation (10 files)
python -m src.eval.evaluate_all_methods \
    --checkpoint checkpoints/masknet_best.pth \
    --max-files 10 \
    --no-plot

# 3. Run demo
python demo.py --checkpoint checkpoints/masknet_best.pth
\end{lstlisting}

\subsection{Most Used Commands}

\begin{lstlisting}[language=bash]
# Train (standard configuration)
python -m src.train.train_mask_model

# Evaluate all methods
python -m src.eval.evaluate_all_methods \
    --checkpoint checkpoints/masknet_best.pth --plot

# Run demo
python demo.py --checkpoint checkpoints/masknet_best.pth

# Check checkpoint
python -c "import torch; ckpt = torch.load(
    'checkpoints/masknet_best.pth', map_location='cpu', 
    weights_only=False); print(ckpt['metrics'])"

# View logs
Get-Content checkpoints/train_*/training.log -Wait -Tail 30
\end{lstlisting}

\section{PowerShell Specific Commands}

\subsection{File Operations}

\begin{lstlisting}[language=powershell]
# List files with size
Get-ChildItem data/librispeech_demand/noisy/train | 
    Select-Object Name, 
        @{Name='MB';Expression={[math]::Round($_.Length/1MB,2)}}

# Find large files
Get-ChildItem -Recurse | 
    Where-Object {$_.Length -gt 100MB} | 
    Sort-Object Length -Descending

# Count files by extension
Get-ChildItem -Recurse | 
    Group-Object Extension | 
    Select-Object Name, Count

# Get directory size
Get-ChildItem data/librispeech_demand/noisy -Recurse | 
    Measure-Object -Property Length -Sum

# Copy with progress
Copy-Item -Path source/* -Destination dest/ -Recurse -Verbose
\end{lstlisting}

\subsection{Process Management}

\begin{lstlisting}[language=powershell]
# Check Python processes
Get-Process python

# Kill Python process
Stop-Process -Name python -Force

# Monitor GPU usage (if nvidia-smi available)
nvidia-smi -l 1

# Check memory usage
Get-Process python | Select-Object Name, 
    @{Name='MemMB';Expression={[math]::Round($_.WS/1MB,2)}}
\end{lstlisting}

\subsection{Log Monitoring}

\begin{lstlisting}[language=powershell]
# Tail log file
Get-Content training.log -Wait -Tail 50

# Search in logs
Select-String -Path "*.log" -Pattern "error|warning"

# Get last N lines
Get-Content training.log -Tail 100

# Monitor multiple logs
Get-Content checkpoints/*/training.log -Wait
\end{lstlisting}

\section{Environment Variables}

\subsection{Set CUDA Device}

\begin{lstlisting}[language=bash]
# PowerShell
$env:CUDA_VISIBLE_DEVICES="0"
python -m src.train.train_mask_model

# Use specific GPU
$env:CUDA_VISIBLE_DEVICES="1"

# Hide GPU (force CPU)
$env:CUDA_VISIBLE_DEVICES=""
\end{lstlisting}

\subsection{Python Path}

\begin{lstlisting}[language=bash]
# Add project root to Python path
$env:PYTHONPATH="d:\proiect licentiat\project_root"

# Verify
python -c "import sys; print(sys.path)"
\end{lstlisting}

\section{Appendix: File Paths}

\subsection{Important Directories}

\begin{itemize}
    \item \textbf{Data:} \texttt{data/librispeech\_demand/}
    \begin{itemize}
        \item Clean: \texttt{clean/\{train,val,test\}/}
        \item Noise: \texttt{noise/\{train,val,test\}/}
        \item Noisy: \texttt{noisy/\{train,val,test\}/}
        \item VAD Labels: \texttt{vad\_labels/\{train,val,test\}/}
    \end{itemize}
    \item \textbf{Checkpoints:} \texttt{checkpoints/}
    \item \textbf{Results:} \texttt{results/}
    \item \textbf{Logs:} \texttt{logs/} and \texttt{checkpoints/train\_*/}
    \item \textbf{Source:} \texttt{src/}
\end{itemize}

\subsection{Key Files}

\begin{itemize}
    \item \textbf{Config:} \texttt{src/config.py}
    \item \textbf{Best Model:} \texttt{checkpoints/masknet\_best.pth}
    \item \textbf{Requirements:} \texttt{requirements.txt}
    \item \textbf{Demo:} \texttt{demo.py}
    \item \textbf{README:} \texttt{README.md}
\end{itemize}

\end{document}
